% =====================================================================
%  Section: Frequency-selective electro-mechanical coupling to the
%           tower side-to-side mode (OpenFAST-FMU + LEOGO)
%  Required in preamble:
%    \usepackage{amsmath, amssymb}
%    \usepackage{graphicx}
%    \usepackage{booktabs}
%    \usepackage{siunitx}
%  Figures produced by:
%    casestudies/dyn_sim/sweep_resonance.py
%    -> casestudies/dyn_sim/results/sweep/resonance_curve.png
%       casestudies/dyn_sim/results/sweep/resonance_summary.csv
%  Adjust the figure path (fig/...) to your own build tree.
% =====================================================================

\section{Frequency-Selective Coupling to the Tower Side-to-Side Mode}
\label{sec:tower-ss-resonance}

The forced-response study of \S\ref{sec:forcedresponse} demonstrated the
electrically driven \emph{drivetrain torsional} resonance. This section turns to
the second mechanical mode of interest for the LEOGO offshore platform: the
first \emph{tower side-to-side} (SS) bending mode. The motivation is the central
research question of the project --- whether a grid disturbance in the platform
network can propagate into, and resonantly excite, a lightly damped mechanical
mode of the wind turbine. The tower SS mode is a natural candidate because its
natural frequency nearly coincides with the LEOGO grid's own electromechanical
oscillation, so any grid oscillation in that band falls exactly where the
mechanical coupling is strongest.

\subsection{Mechanical configuration and frequency coincidence}
\label{subsec:ss-config}

In the OpenFAST (ElastoDyn) model used for the co-simulation only two structural
degrees of freedom are enabled: the generator DOF (\texttt{GenDOF}) and the
first side-to-side tower bending mode (\texttt{TwSSDOF1}). The flapwise,
edgewise, fore-aft-tower and drivetrain-torsion DOFs are all disabled
(\texttt{FlapDOF}, \texttt{EdgeDOF}, \texttt{TwFADOF1}, \texttt{DrTrDOF} =
\texttt{False}). The high-fidelity model therefore \emph{isolates} the tower SS
mode as the single free structural mode driven by the generator, which makes it
the clean target for the interaction study. The turbine operates in Region~2
(hub-height wind \SI{8}{\metre\per\second}, rotor speed
\(\approx\SI{5.70}{rpm}\)), below rated.

The tower SS first mode sits at \(f_{\mathrm{SS}} \approx \SI{0.234}{\hertz}\).
The LEOGO network, dominated by the platform gas-turbine synchronous machines,
has a centre-of-inertia (COI) electromechanical mode at
\(f_{\mathrm{grid}} \approx \SI{0.226}{\hertz}\) --- essentially on top of the
tower mode. This near-coincidence is the physically relevant observation: a
frequency disturbance emitted by the grid in its own electromechanical band
coincides with the frequency at which the tower is most responsive.

\subsection{Torque path and the ROSCO constraint}
\label{subsec:ss-torque-path}

Unlike the reduced model, the electrical\(\to\)mechanical path into the OpenFAST
\emph{structure} is not open by default. ROSCO (\texttt{VSContrl}=5) owns the
generator torque internally and rejects the external torque and power commands
(\texttt{GenSpdOrTrq}, \texttt{ElecPwrCom}) passed across the FMU interface. The
only channel that reaches the OpenFAST structural solver is the ROSCO
\emph{open-loop} generator-torque table (\texttt{OL\_Mode}=1,
\texttt{Ind\_GenTq}=2), read from \texttt{OLInput\_ROSCO.dat} at initialisation.
We therefore prescribe the torque disturbance directly through that table as a
sinusoidal modulation about the Region-2 operating torque
\(T_0 = \SI{11.02}{\mega\newton\metre}\),
\begin{equation}
  T_{\mathrm{gen}}(t) = T_0\,\big[\,1 + a\,
    \sin\!\big(2\pi f\,(t-t_{\mathrm{on}})\big)\,\big]\,
    \mathbb{1}[t \ge t_{\mathrm{on}}],
  \qquad a = 0.10,
  \label{eq:ss-forcing}
\end{equation}
i.e.\ a \(\pm\SI{10}{\percent}\) (\(\pm\SI{1.1}{\mega\newton\metre}\)) generator-torque
oscillation switched on at \(t_{\mathrm{on}}=\SI{10}{\second}\). It must be
stated plainly that the torque source here is \emph{imposed} rather than the
product of a closed grid\(\to\)torque feedback: the ROSCO controller blocks the
genuine electrical path in this FMU build. What the experiment does establish is
genuine OpenFAST structural physics, namely the \emph{frequency selectivity}
with which the tower amplifies a generator-torque oscillation. The
complementary reduced-model study of \S\ref{subsec:fr-simplified}, in which the
electromagnetic-torque path is open by construction, provides the genuinely
closed electrical-to-mechanical demonstration.

This choice of a modal-frequency probe is deliberate. A realistic
grid-frequency transient following a load step is slow and one-sided --- a dip
of order \SI{90}{\milli\hertz} developing over tens of seconds --- so almost all
of its spectral energy lies near DC and virtually none at \SI{0.234}{\hertz}.
Such a disturbance, replayed through a droop law, was found to leave the tower SS
response essentially unchanged. To characterise the coupling itself we therefore
probe the system at, and around, the modal frequency.

\subsection{Sweep method}
\label{subsec:ss-method}

For each probe frequency \(f\) the coupled LEOGO\,+\,OpenFAST-FMU co-simulation
is run for \SI{60}{\second} with a fixed communication step of
\SI{0.01}{\second} and \emph{no} grid load event, so that the sinusoidal
generator torque of \eqref{eq:ss-forcing} is the only disturbance (a clean
single-frequency probe). The tower response is read from the tower-top
side-to-side acceleration \texttt{YawBrTAyp}. To separate the response at the
forcing frequency from the broadband ambient motion, a sinusoid at \(f\) is
fitted by least squares over the settled window \(t \in [25,60]\,\si{\second}\),
\begin{equation}
  y(t) \approx \alpha\sin(2\pi f t) + \beta\cos(2\pi f t) + \gamma,
  \qquad
  A_{\mathrm{SS}} = \sqrt{\alpha^2 + \beta^2},
  \label{eq:ss-fit}
\end{equation}
and the fitted amplitude \(A_{\mathrm{SS}}\) is taken as the response metric.
The sweep, table and figure are produced by
\texttt{casestudies/dyn\_sim/sweep\_resonance.py}.

\subsection{Results}
\label{subsec:ss-results}

The sweep traces out a sharp resonance peaking in the tower SS band
(Table~\ref{tab:ss-sweep}, Figure~\ref{fig:ss-curve}). Far from resonance
(\SIrange{0.10}{0.15}{\hertz}) the fitted side-to-side amplitude is negligible,
\(A_{\mathrm{SS}} \approx \SI{5e-3}{\metre\per\second\squared}\): the imposed
torque oscillation barely moves the tower. Approaching the modal frequency the
response climbs steeply, reaching \(A_{\mathrm{SS}} \approx
\SI{0.29}{\metre\per\second\squared}\) at \SI{0.245}{\hertz} --- an amplification
of about \(60\times\) relative to the off-resonance floor --- and falls off just
as steeply above \SI{0.26}{\hertz}. The peak lies in the
\SIrange{0.234}{0.245}{\hertz} band, consistent with the tower SS first mode; the
measured maximum sits marginally above the nominal \SI{0.234}{\hertz}, an offset
attributable to the coarse frequency grid and to the finite settling within the
measurement window. The narrowness of the curve --- amplitudes above the
half-power level (\(A_{\mathrm{SS}} \gtrsim \SI{0.20}{\metre\per\second\squared}\))
span only \(\approx\SI{0.02}{\hertz}\) --- indicates a lightly damped, high-\(Q\)
mode (\(Q \sim 10\), i.e.\ \(\zeta \sim \SI{5}{\percent}\)).

Crucially, the LEOGO grid electromechanical mode at \SI{0.226}{\hertz} falls on
the steep rising flank of this resonance curve, where the tower response has
already grown to \(A_{\mathrm{SS}} \approx \SI{0.14}{\metre\per\second\squared}\)
--- some \(30\times\) the off-resonance floor. A grid oscillation in the
platform's own electromechanical band therefore lands squarely in the frequency
range to which the tower is strongly susceptible.

\begin{table}[t]
  \centering
  \caption{OpenFAST-FMU frequency sweep of the tower side-to-side response:
  least-squares fitted amplitude of \texttt{YawBrTAyp} at the forcing frequency
  (\(\pm\SI{10}{\percent}\) generator-torque modulation, Region~2). Rows in bold
  mark the LEOGO grid mode and the tower SS mode.}
  \label{tab:ss-sweep}
  \begin{tabular}{lc}
    \toprule
    \(f\) [\si{\hertz}] &
      \(A_{\mathrm{SS}}\) [\si{\metre\per\second\squared}] \\
    \midrule
    0.100          & 0.0047 \\
    0.150          & 0.0071 \\
    0.190          & 0.0473 \\
    0.210          & 0.0680 \\
    \textbf{0.226} & \textbf{0.1399} \\
    \textbf{0.234} & \textbf{0.2591} \\
    0.245          & 0.2873 \\
    0.260          & 0.0734 \\
    0.300          & 0.0296 \\
    0.400          & 0.0195 \\
    \bottomrule
  \end{tabular}
\end{table}

\begin{figure}[t]
  \centering
  % \includegraphics[width=0.8\linewidth]{fig/resonance_curve.png}
  \caption{OpenFAST-FMU model: fitted tower-top side-to-side acceleration
  amplitude (\texttt{YawBrTAyp}) versus generator-torque modulation frequency
  (\(\pm\SI{10}{\percent}\), Region~2). The response peaks sharply in the tower
  side-to-side band (dashed line, \SI{0.234}{\hertz}); the LEOGO grid
  electromechanical mode (dotted line, \SI{0.226}{\hertz}) sits on the rising
  flank of the resonance.}
  \label{fig:ss-curve}
\end{figure}

\subsection{Discussion}
\label{subsec:ss-discussion}

The sweep is direct evidence of frequency-selective electro-mechanical coupling:
the OpenFAST tower structure amplifies a generator-torque oscillation by roughly
\(60\times\) when the oscillation frequency matches its side-to-side eigenmode,
and hardly at all otherwise. Combined with the near-coincidence of the LEOGO grid
electromechanical mode (\SI{0.226}{\hertz}) and the tower SS mode
(\SI{0.234}{\hertz}), this makes the grid\(\to\)mechanical propagation
physically plausible rather than merely postulated.

Two caveats frame the interpretation. First, as noted in
\S\ref{subsec:ss-torque-path}, the torque is injected through the ROSCO
open-loop table because the controller blocks the genuine grid\(\to\)torque path
in this FMU build; the resonant \emph{amplification} is real structural physics,
but the excitation is prescribed. Second, the result is best read as a
\emph{susceptibility} (worst-case) statement: the LEOGO COI mode is strongly
damped and does not ring spontaneously, so a sustained tower oscillation would
require a persistent forcing in the \SIrange{0.22}{0.25}{\hertz} band (for
example from poorly damped operating conditions or sub-synchronous control
interactions). The finding is that \emph{if} the grid oscillates in this band,
the frequency coincidence makes the coupling to the tower strong. The genuinely
open electrical-to-mechanical demonstration, without any imposed torque, is
provided by the reduced-model torsional study of \S\ref{sec:forcedresponse};
the two models are complementary --- the reduced model proves the open coupling
path, and the high-fidelity model quantifies the frequency selectivity of a
realistic tower mode.
